\cotefile{100}
\scanpage{15}
montrer que la donnée de descente envisagée est effective. Or on \deleted{se ramène facilement au cas où $S$ est local, pt fermé $s$, $U=S-\{s\}$. Par factorisation (\unsure{abandonnant|abandonne|où utilisant} \deleted{l'hypothèse} \inserted{le fait} que $S'$ normal) on est ramené à un morphisme $S'\to S$ qui est un épimorphisme effectif. \deleted{\unsure{donc}} (\unsure{$\overset{\text{''}}{\operatorname{Spec}} K'\to\overset{\text{''}}{\operatorname{Spec}}K$|$\overset{\text{''}}{\operatorname{Spec}} k'\to\overset{\text{''}}{\operatorname{Spec}}k$}) \deleted{$R'/R$ de \unsure{rang~1 sur $R$|deg. 1 sur $R$} (?)}} \margin{\deleted{\unsure{question} effectivité des descentes à régler}} \deleted{utilisant l'effectivité des descentes en vertu de FGA on est ramené à montrer qu'il \deleted{\unsure{Il suffit alors que}} existe sur $A'$ une polarisation stable par la donnée de descente} \unsure{a} un th. de descente qui nous assure cela, cf plus bas.
% archive sheet number: 14

\margin{\uline{?!}}
e) Cas général. \deleted{Le cas} On a déjà vu qu'on peut supposer $S$ réduit. Le cas $S$ \uline{local} relève de d).
% slanted
\margin{\uline{\unsure{Lemmes|Lemme}} (\unsure{indépendant|indépendt} de l'hyp. faite \unsure{dans|ds} la partie $2^{\text{ème}}$ du th)\,: pour que $(A_U,M,\varphi)$ \unsure{appartienne} à l'image essentielle, il f. et s. que, pour tt pt \deleted{\unsure{fermé}} $t$ de $T$, il y ait une solution sur $\operatorname{Spec}\mathcal O_{S,t}$. Il suffit \unsure{même|=} de se borner aux $t$ \uline{fermés}, ou au choix aux $t$ \unsure{maximaux} \unsure{dans|des} ``tels que \unsure{les $s$|des $s$}'' qu'il y a une sol. \unsure{loc.} dans $\operatorname{Spec}\mathcal O_{S,t}-\{t\}$.}
Cela permet déjà, pour tt $s\in S$, de trouver un schéma abélien $A_{(s)}$ sur $\operatorname{Spec}\mathcal O_{S,s}=S_{\unsure{\uparrow|1}}^{(s)}$, qui soit une solution au Pb induit sur $S_{(s)}$ (avec $U_{(s)}=U\times_S S_{(s)}$, $T_{(s)}=T\cap S_{(s)}$ etc $\ldots$). Par unicité, les $A_{(s)}$ forment un ``système local de schémas abéliens'', i.e. pour $t$ gén. de $s$, on a un isom can $A_{(\deleted{\unsure{s}}t)}\simeq A_{(s)}\times_{S_{(s)}}S_{(t)}$, avec transitivité. \deleted{Supposons que \unsure{pour}} Or on a un foncteur \uline{pl. fid.} \unsure{Mord|Mod} de la catégorie des schémas de t.f. sur $S$ dans la cat. des systèmes locaux de schémas de \unsure{t.f.|s.f.} sur les $S_{(s)}$. Supposons qu'on ait prouvé que $\{A_{(s)}\}_s$ est dans l'image essentielle. Alors on a donc un schéma \deleted{abél.} $A$ sur $S$, néc. abélien $\ldots$ (par la pleine fidélité précédente). On a

